\section{The scaling toolbox}\label{sec:toolbox}

Throughout this section $U\in\R^{n\times d}$ is Parseval, $P=UU^T$,
$p=\diag P$, and $L=L_P=\Diag(p)-P\circ P$.

\subsection{Projections, distances, and the squared-Gram Laplacian}

\begin{lemma}[Frame versus projection distance]\label{lem:align}
Let $U,W\in\R^{n\times d}$ be Parseval with projections $P,Q$, and let
$\sigma_1,\dots,\sigma_d\in[0,1]$ be the singular values of $U^TW$. Then
\[
 \min_{R\in O(d)}\fro{U-WR}^2=2\sum_k(1-\sigma_k)\le
 \fro{P-Q}^2=2\sum_k(1-\sigma_k^2)=2\fro{(I-P)W}^2\le2\fro{U-W}^2 .
\]
Right multiplication by $R\in O(d)$ preserves Parsevalness and every row norm.
\end{lemma}

\begin{proof}
$\fro{U-WR}^2=2d-2\tr(U^TWR)$, maximised over $R$ by the polar factor, giving
$2d-2\sum_k\sigma_k$. Next $\fro{P-Q}^2=2d-2\tr(PQ)=2d-2\fro{U^TW}^2$ and
$\fro{(I-P)W}^2=d-\fro{U^TW}^2$. Finally $1-\sigma\le1-\sigma^2\le2(1-\sigma)$
for $\sigma\in[0,1]$.
\end{proof}

\begin{lemma}[Energy identity]\label{lem:energy}
$L$ is the Laplacian of the weighted graph with weights $P_{ij}^2$ $(i\ne j)$,
$L\succeq0$, $L\one=0$, $L_{I-P}=L_P$, and for $x\in\R^n$, $X=\Diag(x)$,
\[
 x^TLx=\tfrac12\sum_{i,j}P_{ij}^2(x_i-x_j)^2=\fro{(I-P)XU}^2
 =\tfrac12\fro{[X,P]}^2 .
\]
\end{lemma}

\begin{proof}
$\sum_jP_{ij}^2=(P^2)_{ii}=p_i$, so the rows of $\Diag(p)-P\circ P$ sum to zero
and its off-diagonal entries are $-P_{ij}^2$; this gives the first formula. Next
$\fro{(I-P)XU}^2=\tr(XPX(I-P))=\sum_ix_i^2p_i-\sum_{i,j}x_ix_jP_{ij}^2$, and
$[X,P]_{ij}=(x_i-x_j)P_{ij}$. For $I-P$ the off-diagonal weights are again
$P_{ij}^2$.
\end{proof}

\begin{lemma}[Scaled projections]\label{lem:scaled}
For $w\in\R^n_{>0}$ put $D=\Diag(w)$, $M=(U^TD^2U)^{-1/2}$ and $W=DUM$. Then
$W$ is Parseval with column space $D\im U$, $\opn M\le1/\min_iw_i$, and
\[
 \fro{WW^T-P}^2=2\fro{(I-P)W}^2\le\frac{2\,w^TLw}{\min_iw_i^2} .
\]
\end{lemma}

\begin{proof}
$W^TW=M U^TD^2U M=I$ and $U^TD^2U\succeq(\min w_i)^2I$. By \cref{lem:align}
and \cref{lem:energy},
$\fro{(I-P)W}^2\le\opn M^2\fro{(I-P)DU}^2=\opn M^2\,w^TLw$.
\end{proof}

\subsection{The log-determinant potential}

For $s\in\R^n$ let $D_s=\Diag(e^{s})$, let $P(s)$ be the projection onto
$D_s\im U$, and
\[
 F(s)=\tfrac12\log\det(U^TD_s^2U).
\]

\begin{lemma}\label{lem:potential}
$F$ is smooth and convex, $F(s+c\one)=F(s)+cd$, and
\[
 \nabla F(s)=\diag P(s),\qquad \nabla^2F(s)=2L_{P(s)} .
\]
Moreover, if $\norm{s}_\infty\le r$ then
\begin{equation}\label{eq:hessian-comparison}
 e^{-4r}L_P\preceq L_{P(s)}\preceq e^{4r}L_P .
\end{equation}
\end{lemma}

\begin{proof}
Write $G=U^TD_s^2U=\sum_ie^{2s_i}u_iu_i^T$. Then
$\partial_iF=e^{2s_i}u_i^TG^{-1}u_i=P(s)_{ii}$ because
$P(s)=D_sUG^{-1}U^TD_s$, and differentiating once more gives
$\partial_j\partial_iF=2\delta_{ij}P(s)_{ii}-2P(s)_{ij}^2$, i.e.\
$\nabla^2F=2L_{P(s)}\succeq0$. For~\eqref{eq:hessian-comparison} let
$W_s=D_sUM_s$ with $M_s=G^{-1/2}$, so $\opn{M_s}\le e^r$, $\opn{D_s}\le e^r$.
Since $(I-P(s))D_sP=(I-P(s))W_sM_s^{-1}U^T=0$,
\[
 (I-P(s))\Diag(x)W_s=(I-P(s))D_s(I-P)\Diag(x)UM_s ,
\]
and \cref{lem:energy} for $W_s$ gives $x^TL_{P(s)}x\le e^{4r}x^TL_Px$. The
lower bound is the upper bound applied to the Parseval frame $W_s$ and the
scaling $-s$, since $D_{-s}W_s=UM_s$ has column space $\im U$.
\end{proof}

\begin{lemma}[Exact scaling inequalities]\label{lem:scaling-ineq}
For $z\in\R^n_{>0}$ let $r$ be the diagonal of the projection onto
$\Diag(\sqrt z)\im U$. Then, coordinatewise,
\[
 Lz\le z\circ(r-p),\qquad L(z^{-1})\le-z^{-1}\circ(r-p).
\]
\end{lemma}

\begin{proof}
Let $G=U^T\Diag(z)U$. The matrix $(G-z_iI)G^{-1}(G-z_iI)$ is positive
semidefinite; evaluating its quadratic form at $u_i$ and using
$u_i^TGu_i=\sum_jP_{ij}^2z_j$, $\norm{u_i}^2=p_i$, $z_iu_i^TG^{-1}u_i=r_i$ gives
$\sum_jP_{ij}^2z_j-2z_ip_i+z_ir_i\ge0$, which is the first inequality. Let
$V\in\R^{n\times(n-d)}$ be an isometry onto $(\im U)^\perp$. The subspaces
$\Diag(\sqrt z)\im U$ and $\Diag(z^{-1/2})\im V$ are orthogonal complements,
so the projection onto the latter has diagonal $\one-r$, while $VV^T=I-P$ has
diagonal $\one-p$ and Laplacian $L$ (\cref{lem:energy}). The first inequality
for $(V,z^{-1})$ is the second inequality.
\end{proof}

\subsection{Static balancing}

\begin{definition}\label{def:poisson}
A weighted graph Laplacian $L$ on $n$ vertices has \emph{Poisson constant}
$K$ if for every $f\perp\one$ there is $g$ with $Lg=f$ and
$\norm{g}_\infty\le K\norm f_\infty$.
\end{definition}

If $L$ has a finite Poisson constant then the graph is connected (on each
component the image of $L$ sums to zero). The Poisson constant is the
$\ell_\infty\to\ell_\infty$ norm of the pseudo-inverse modulo constants; it is
the quantity that controls a scaling in the box $\norm s_\infty\le1$.

\begin{lemma}[Median barrier]\label{lem:median}
Let $L$ have Poisson constant $K$, let $\beta\ge0$ with $2\beta K<1$, let
$z\in\R^n_{>0}$, let $m$ be a median of $z$ (so that $\{z\le m\}$ and
$\{z\ge m\}$ both have at least $n/2$ elements), and suppose
\[
 (Lz)_i\le\beta z_i\ \text{ when }z_i>m,\qquad
 (Lz^{-1})_i\le\beta z_i^{-1}\ \text{ when }z_i<m .
\]
Then $\max_iz_i/\min_iz_i\le(1-2\beta K)^{-2}$.
\end{lemma}

\begin{proof}
\emph{Claim.} If $|S|\ge n/2$, $y\ge0$, $y\le m$ on $S$ and $(Ly)_i\le\beta y_i$
for $i\notin S$, then $(1-2\beta K)\max y\le m$. Indeed, $f=\one-\frac
n{|S|}\one_S$ has $f\perp\one$ and $\norm f_\infty\le1$; take $Lg=f$ with
$\norm g_\infty\le K$ and $h=g+K\one\in[0,2K]^n$, so $(Lh)_i=1$ for $i\notin
S$. With $M=\max y$, the vector $v=y-m\one-\beta Mh$ is $\le0$ on $S$ and
satisfies $(Lv)_i\le\beta y_i-\beta M\le0$ off $S$. If $\max v>0$, it is attained
at some $i\notin S$, and $(Lv)_i=\sum_jw_{ij}(v_i-v_j)\le0$ forces $v_j=v_i$
for all neighbours $j$; by connectedness the maximum propagates to $S$, a
contradiction. Hence $y\le m+2\beta KM$, and evaluating at the maximum proves
the claim.

Apply the claim to $y=z$, $S=\{z\le m\}$, and to $y=z^{-1}$, $S=\{z\ge m\}$
(with median $m^{-1}$). This gives $(1-2\beta K)\max z\le m$ and
$(1-2\beta K)m\le\min z$.
\end{proof}

\begin{theorem}[Static balancing]\label{thm:balancing}
Suppose $L_P$ has Poisson constant $K$, and $q\in\R^n$ satisfies
$\sum_iq_i=d$, $\norm{q-p}_\infty\le\beta$ and $\beta K\le\frac14$. Then there
is a Parseval $W$ with $\diag(WW^T)=q$ and
\[
 \fro{U-W}^2\le\fro{UU^T-WW^T}^2\le8\norm{q-p}_1 .
\]
$W$ is obtained from $U$ by a row scaling with ratio of scales at most $2$,
whitening, and an orthogonal right factor.
\end{theorem}

\begin{proof}
Let $\Phi(s)=F(s)-\ip qs$, so $\partial_i\Phi=P(s)_{ii}-q_i$, and let $s^*$
minimise $\Phi$ over the cube $[0,1]^n$. Write $r=\diag P(s^*)$. Coordinatewise
optimality gives
\begin{equation}\label{eq:kkt}
 r_i\le q_i\ \text{ if }s^*_i>0,\qquad r_i\ge q_i\ \text{ if }s^*_i<1 .
\end{equation}
Let $z=e^{2s^*}$ with median $m$; then $1\le\min z\le m\le\max z\le e^2$. If
$z_i>m$ then $s_i^*>0$, so by
\cref{lem:scaling-ineq} and~\eqref{eq:kkt}
$(Lz)_i\le z_i(r_i-p_i)\le z_i(q_i-p_i)\le\beta z_i$; if $z_i<m$ then $s_i^*<1$
and likewise $(Lz^{-1})_i\le\beta z_i^{-1}$. By \cref{lem:median},
$\max s^*-\min s^*\le-\log(1-2\beta K)\le\log2<1$. Hence $s^*$ does not meet
both faces: either $s^*_i>0$ for all $i$, and then $r\le q$, or some $s^*_j=0$,
all $s_i^*<1$, and then $r\ge q$. Since $\sum r_i=d=\sum q_i$, in both cases
$r=q$.

Let $w=e^{s^*}$, $D=\Diag(w)$, $m_0=\min w$, $W_0=DU(U^TD^2U)^{-1/2}$, and
$\delta=q-p$. \Cref{lem:scaling-ineq} with $r=q$ gives $Lz\le z\circ\delta$
for $z=w^2$, hence $z^TLz\le\sum_iz_i^2\delta_i\le16m_0^4\norm\delta_1$, using
$z_i\le4m_0^2$. Since $(z_i-z_j)^2\ge4m_0^2(w_i-w_j)^2$,
\cref{lem:energy} gives $w^TLw\le4m_0^2\norm\delta_1$, and
\cref{lem:scaled} gives $\fro{W_0W_0^T-P}^2\le8\norm\delta_1$. Finally take
$W=W_0R$ with $R$ from \cref{lem:align}.
\end{proof}

\begin{remark}
No existence theorem for scalings, no flow, and no genericity is used: the box
supplies a minimiser, the median barrier keeps it away from one face, and the
trace identity $\sum_ir_i=d$ does the rest.
\end{remark}

\subsection{Local correction in the dual norm}

The balancing theorem needs $\ell_\infty$ control of the target. The next lemma
handles a different kind of target error: one that is small in the norm dual
to a Laplacian with a spectral gap, although possibly large in $\ell_\infty$.

\begin{lemma}[Local correction]\label{lem:local}
Let $W$ be Parseval with projection $P_W$ and diagonal $p_W$. Let $J$ be
symmetric with $J\succeq\sigma I$ on $\one^\perp$ ($\sigma>0$) and
\[
 c_1J\preceq L_{P_W}\preceq C_1J\quad\text{on }\one^\perp .
\]
Let $f\perp\one$ satisfy $\ip fx^2\le E\,x^TJx$ for $x\perp\one$, and let
$0<R\le r$ with
\begin{equation}\label{eq:local-small}
 E<\gamma^2\sigma R^2,\qquad \gamma:=2c_1e^{-4r}.
\end{equation}
Then there is $s\perp\one$ with $\norm s_\infty<R$, $s^TJs\le E/\gamma^2$, such
that the Parseval frame $W'$ obtained from $\Diag(e^s)W$ by whitening satisfies
$\diag(W'W'^T)=p_W+f$ and $L_{P_{W'}}\succeq c_1e^{-4r}J$ on $\one^\perp$.
After an orthogonal right factor,
\[
 \fro{W'-W}^2\le\fro{P_{W'}-P_W}^2\le2C_1e^{4r}E/\gamma^2 .
\]
\end{lemma}

\begin{proof}
Let $F$ be the potential of $W$ and $\Phi(s)=F(s)-F(0)-\ip{p_W+f}s$; minimise
$\Phi$ over the compact convex set $\Omega=\{s\perp\one:\norm s_\infty\le R\}$.
For $s\in\Omega$ with $\norm s_\infty=R$, \cref{lem:potential} gives
\[
 \ip{\nabla\Phi(s)}s=\int_0^12s^TL_{P(us)}s\,du-\ip fs
 \ge2c_1e^{-4r}s^TJs-\sqrt{E\,s^TJs}>0,
\]
because $s^TJs\ge\sigma\norm s_2^2\ge\sigma R^2>E/\gamma^2$. So $-s$ is a
feasible descent direction and the minimiser $s$ lies in the relative interior
of $\Omega$. There $\nabla\Phi(s)\in\R\one$, and $\sum_i\nabla\Phi(s)_i=d-d-0=0$,
so $\diag P(s)=p_W+f$. Pairing $\nabla\Phi(s)=0$ with $s$ as above gives
$\gamma\, s^TJs\le\sqrt{E s^TJs}$, i.e.\ $s^TJs\le E/\gamma^2$, and then
$\norm s_\infty\le\norm s_2<R$ by~\eqref{eq:local-small}. The Laplacian bound
is~\eqref{eq:hessian-comparison}. For the distance, \cref{lem:scaled} with
$w=e^s$, $\min w_i\ge e^{-r}$ and $|e^{s_i}-e^{s_j}|\le e^r|s_i-s_j|$ gives
\[
 \fro{P_{W'}-P_W}^2\le2e^{2r}(e^s)^TL_{P_W}e^s\le2e^{4r}s^TL_{P_W}s
 \le2C_1e^{4r}s^TJs ,
\]
and \cref{lem:align} supplies the right factor.
\end{proof}

\subsection{Poisson constants from a dense core}

\begin{lemma}[Dense core with exceptional vertices]\label{lem:core}
Let $L$ be the Laplacian of nonnegative symmetric weights $w_{ij}$ ($w_{ii}=0$)
on $G\sqcup B$, $G\neq\varnothing$, $|B|=b$. Assume
\[
 \sum_{j\in G}\min\{w_{ij},w_{kj}\}\ge\rho>0\quad(i,k\in G),\qquad
 \max_i\sum_jw_{ij}\le D_{\max},
\]
and, if $B\ne\varnothing$, that $C=L_{BB}$ is invertible with
$\norm{C^{-1}}_{\infty\to\infty}\le R$. Then $L$ has Poisson constant
\[
 K=\frac{1+\min\{b,D_{\max}R\}}{\rho}+R
\]
(with $R=b=0$ if $B=\varnothing$).
\end{lemma}

\begin{proof}
Write $L=\begin{pmatrix}A&-E\\-E^T&C\end{pmatrix}$ with $E=(w_{ij})_{i\in
G,j\in B}\ge0$. $C$ is a symmetric $Z$-matrix with $C\one_B=E^T\one_G\ge0$;
since it is invertible and positive semidefinite it is positive definite, and
$C^{-1}\ge0$ entrywise (if $Cx=y\ge0$ then, with $x_-$ the negative part,
$0\le x_-^Ty=x_-^TCx\le-x_-^TCx_-$, so $x_-=0$). Put $H=C^{-1}E^T\ge0$ and
$T=EC^{-1}=H^T$; then $H\one_G=\one_B$ and $\one_G^TT=\one_B^T$. Each row sum of
$T$ is at most $b$ (entries of a column of $T$ are bounded by its column sum
$1$) and at most $D_{\max}R$ (as $T\one_B\le R\,E\one_B$).

The Schur complement $S=A-EC^{-1}E^T$ is the
Laplacian on $G$ with weights $\tilde w_{ik}=w_{ik}+(EC^{-1}E^T)_{ik}\ge w_{ik}$
($i\ne k$)
(its rows sum to zero because $C^{-1}E^T\one_G=\one_B$). Let $f\perp\one$ with
$\norm f_\infty\le1$, and put $h=f_G+Tf_B$; then $\sum h=\sum f=0$ and
$\norm h_\infty\le\Phi:=1+\min\{b,D_{\max}R\}$. The weights $\tilde w$ inherit
the overlap bound, so the graph on $G$ is connected and $Sy=h$ is solvable. If
$i,k\in G$ maximise and minimise $y$,
\[
 h_i-h_k=\sum_{j\in G}\bigl[\tilde w_{ij}(y_i-y_j)+\tilde w_{kj}(y_j-y_k)\bigr]
 \ge\rho\,(y_i-y_k),
\]
so after adding a constant $\norm y_\infty\le\Phi/\rho$. Then $g_G=y$,
$g_B=Hy+C^{-1}f_B$ solves $Lg=f$ (block substitution), and
$\norm{g_B}_\infty\le\Phi/\rho+R$ because $H$ is row-stochastic.
\end{proof}

\begin{corollary}\label{cor:core}
Suppose $b\le n/10$, $\gamma>0$, and every $i\in G$ has at least $\frac45n$
indices $j\ne i$ with $w_{ij}\ge\gamma/n$. Then the overlap hypothesis holds with
$\rho=\gamma/2$. Moreover:
\begin{enumerate}[label=(\alph*)]
\item if $L\succeq\lambda\Pi$ with $\lambda>0$, then $R\le2b/\lambda$;
\item if $w_{ij}=P_{ij}^2$ for a projection $P$, $a>0$, $P_{ii}\ge a/2$ on $B$,
$\sum_{i\in B}P_{ii}\le\frac14$, and $\max_iP_{ii}\le C_2a$, then
$R\le4/a$ and $D_{\max}R\le4C_2$.
\end{enumerate}
\end{corollary}

\begin{proof}
Two sets of at least $\frac45n$ indices share at least $\frac35n$, of which at
least $\frac12n$ lie in $G$, giving overlap $\ge\gamma/2$. (a) For $x$ supported
on $B$, $x^TCx=x^TLx\ge\lambda(\norm x^2-(\sum x_i)^2/n)\ge\lambda(1-b/n)\norm
x^2$, so $|(C^{-1})_{ij}|\le\opn{C^{-1}}\le(\lambda(1-b/n))^{-1}$; sum over a
row. (b) By $P_{ij}^2\le P_{ii}P_{jj}$,
$(C\one_B)_i=P_{ii}-\sum_{j\in B}P_{ij}^2\ge P_{ii}(1-\sum_{j\in B}P_{jj})\ge
a/4$; so $C$ is strictly diagonally dominant, $C^{-1}\ge0$, and $C^{-1}\one_B\le
(4/a)\one_B$. Also $D_{\max}\le\max_iP_{ii}$.
\end{proof}

\subsection{From a seed to an exact frame}

The two constructions produce the following object.

\begin{definition}[Seed]\label{def:seed}
Let $\Theta\ge2$ and $t>0$. A \emph{$(\Theta,t)$-seed} for an $n\times d$ frame
$X$ is a Parseval $V\in\R^{n\times d}$ with
\[
 \fro{V-X}^2\le\Theta t^2d,\qquad
 L_{VV^T}\text{ has Poisson constant }K\le\frac{\Theta}{at^2},\qquad
 \max_i\bigl|(VV^T)_{ii}-a\bigr|\le\frac{at^2}{4\Theta}.
\]
\end{definition}

\begin{corollary}\label{cor:seed}
If $X$ has a $(\Theta,t)$-seed, there is an ENP frame $W$ with
$\fro{X-W}^2\le4\Theta t^2d$.
\end{corollary}

\begin{proof}
Apply \cref{thm:balancing} to $V$ with $q=a\one$:
$\beta K\le\frac14$, and $\fro{V-W}^2\le8n\cdot\frac{at^2}{4\Theta}\le2t^2d$.
Then $\fro{X-W}^2\le2\fro{X-V}^2+2\fro{V-W}^2\le(2\Theta+4)t^2d\le4\Theta
t^2d$.
\end{proof}
